% GPT-6 Astra Pro; inherited effort, exact setting not exposed; /root.
\section[HD-F1-T17: symplectic-type multiplicities at the original source]{HD-F1-T17: original inertia excludes symplectic-type multiplicities through six}
\label{sec:followup-symplectic-multiplicity}

Retain the original faithful absolutely irreducible coefficient
representation $G\subset\operatorname{GL}_K(V)$ of the remaining
T2-OTHER gate. Thus $m=\dim_K V\ge27$ and the actual generating
inertia tuple satisfies
\[
 g_1\cdots g_b=1,\qquad \langle g_1,\ldots,g_b\rangle=G,
 \quad b\ge12,\quad g_j\ne1,\qquad
 \sum_{j=1}^b\operatorname{rank}(g_j-1)=2m+6.
 \tag{F1T17.1}
\]
Consider an actual normal subgroup $R\lhd G$ whose centre $Z_R$
is cyclic and acts by scalars on $V$, such that
$R/Z_R$ is a nonzero elementary abelian $p$-group and the
commutator pairing
\[
 R/Z_R\times R/Z_R\longrightarrow Z_R[p]
 \tag{F1T17.2}
\]
is nondegenerate. Then the restriction of $V_{\mathbf C}$ to
$R$ is isotypic. Write its multiplicity as $k$. No remaining
source has $k\le6$. At $k=7$ its original dimension is one of
$28,56,112$. The centre-scalar hypothesis and the specified normal
subgroup are part of this statement; no classification of the other
normal subgroups is presumed.

\paragraph{1. The actual isotypic evaluation map.}
Let $\xi:Z_R\to\mathbf C^\times$ be the central character.
It is faithful because the original coefficient representation
is faithful. For $h\in R\setminus Z_R$, nondegeneracy gives
$e\in R$ with $ehe^{-1}=zh$ for a nonidentity
$z\in Z_R[p]$. Every irreducible constituent $W$ of
$V_{\mathbf C}|_R$ has central character $\xi$; its character
therefore satisfies
\[
 \chi_W(h)=\xi(z)\chi_W(h),\qquad \chi_W(h)=0.
 \tag{F1T17.3}
\]
The two characters of any two such constituents have positive
inner product
\[
 \frac1{|R|}\sum_{h\in R}
       \overline{\chi_{W_1}(h)}\chi_{W_2}(h)
       =\frac{|Z_R|}{|R|}\dim W_1\dim W_2>0.
 \tag{F1T17.4}
\]
Here the left side is exactly the dimension of
$\operatorname{Hom}_R(W_1,W_2)$: average the action on the actual Hom space,
\[
             f\longmapsto\rho_{W_2}(h)f\rho_{W_1}(h)^{-1}.
\] Its average is the projection onto invariant
homomorphisms and its trace is the displayed sum. A nonzero
homomorphism between irreducible modules is an isomorphism,
by its invariant kernel and image. Thus all constituents are
one actual $W$. Complete reducibility follows directly by
averaging a positive Hermitian form over the finite group.
Put $d=\dim W$ and $M=\operatorname{Hom}_R(W,V_{\mathbf C})$.
The evaluation map is the actual isomorphism
\[
 W\otimes_{\mathbf C}M\xrightarrow{\ \epsilon\ }V_{\mathbf C},
          \qquad w\otimes f\longmapsto f(w),\qquad m=dk.
 \tag{F1T17.5}
\]
Schur's lemma here follows from an eigenvalue of an endomorphism
and its invariant kernel. Consequently the centralizer of $R$
acts under (F1T17.5) by $1_W\otimes T$. This is a matrix
identity on the original complexified coefficient space.
Applying the same averaging formula to $W_1=W_2=W$ now gives
$d^2=|R/Z_R|$. In the eight remaining coefficient fields of T2,
the scalar roots of unity have order four or two, so the prime
in (F1T17.2) is $p=2$. A nondegenerate alternating form on
$R/Z_R$ gives $|R/Z_R|=2^{2a}$ with $a\ge1$, and therefore
$d=2^a$.

\paragraph{2. The ranks inside and outside the actual centralizer.}
Set $C=C_G(R)$. It is normal in $G$, since $R$ is normal.
Every nonidentity $g\in C$ has $T\ne1_M$, and hence
\[
           \operatorname{rank}_V(g-1)
               =d\operatorname{rank}_M(T-1)\ge m/k.
 \tag{F1T17.6}
\]
For $g\notin C$ we prove the uniform bound
\[
                         \operatorname{rank}_V(g-1)\ge m/4.
 \tag{F1T17.7}
\]
First, every $h\in R\setminus Z_R$ has $p$ eigenspaces
of equal dimension $m/p$. Indeed $h^p\in Z_R$, and
conjugation by a suitable $e\in R$ multiplies $h$ by
a primitive $p$th root of unity. It permutes all roots
of the scalar polynomial $X^p-h^p$ cyclically. Thus
\[
          \operatorname{rank}(h-t)\ge m(1-1/p)
                     \quad\hbox{for every scalar }t.
 \tag{F1T17.8}
\]
If some $e\in R$ makes $h=[g,e]\notin Z_R$, the exact
identity from T15 gives
\[
 h-1=((g-1)e-e(g-1))g^{-1}e^{-1},\qquad
 m(1-1/p)\le\operatorname{rank}(h-1)
                        \le2\operatorname{rank}(g-1).
 \tag{F1T17.9}
\]
This proves (F1T17.7) in that case.

In the remaining case every $[g,e]$ lies in $Z_R$.
As $g$ fixes this scalar centre pointwise and $e^p\in Z_R$,
the map $eZ_R\mapsto[g,e]$ is a character from $R/Z_R$
to $Z_R[p]$. The perfect pairing (F1T17.2) gives an
actual $u\in R$ with $[u,e]=[g,e]$ for every $e$.
Then $u^{-1}g\in C$, and under (F1T17.5) one has
\[
                       g=\rho_W(u)\otimes T.
 \tag{F1T17.10}
\]
Since $g\notin C$, this $u$ is noncentral. Its $p$
eigenspaces on $W$ have dimension $d/p$. The finite-order
matrix $g$ is semisimple, and so is $T$: raising (F1T17.10)
to a common multiple of the orders of $u$ and $g$ gives
a finite-order equation for $T$. For each eigenvalue $t$
of $T$, at most $d/p$ eigenvectors in $W$ can make the
tensor eigenvalue one. Summing their exact multiplicities
gives $\dim V^g\le dk/p=m/p$ and in particular (F1T17.7).
This calculation does not replace $u^{-1}g$ by a scalar.

\paragraph{3. Product one forces at least three outside-centralizer terms.}
The quotient $G/C$ contains the image of $R$, whose kernel
on $R$ is exactly $Z_R$. Hence it contains the noncyclic
group $R/Z_R$. That group is noncyclic because its
nondegenerate alternating pairing has nonzero dimension;
a one-dimensional alternating space has zero pairing.
The original tuple in (F1T17.1) projects to a generating
product-one tuple of $G/C$. It has at least three
nonidentity entries. With zero or one entry the product
condition gives the trivial group; with two entries they
are inverse and generate a cyclic group, whereas $G/C$
contains the noncyclic subgroup just exhibited.
Thus at least three original inertia matrices lie
outside $C$.

For $k\le4$, (F1T17.6)--(F1T17.7) give
$2m+6\ge bm/4\ge3m$, so $m\le6$.
For $k\ge4$, the outside-centralizer lower bound is at
least the inside-centralizer bound. Using at least three
outside entries and at least twelve total entries gives
\[
                    2m+6\ge \frac{3m}{4}+\frac{9m}{k}.
 \tag{F1T17.11}
\]
In particular the exact consequences are
\[
 \begin{array}{c|c}
 k&\hbox{bound on the original }m\\ \hline
 1,2,3,4&m\le6\\
 5&11m\le120\\
 6&m\le24\\
 7&m\le168.
 \end{array}
 \tag{F1T17.12}
\]
The first six multiplicities contradict $m\ge27$. For the seventh,
$m=7\cdot2^a$ and $27\le m\le168$ give exactly
$m\in\{28,56,112\}$.
There is no claim that an object or cover exists at any
integer left by this inequality.

\paragraph{Original geometric maps and framework.}
The only tensor isomorphism used is the evaluation map
(F1T17.5) on $V_{\mathbf C}$. All matrix ranks remain
those of the original representation over $K$, since
extension to $\mathbf C$ preserves rank. The quotient
$G\to G/C_G(R)$ is used solely to count the original
nonidentity inertia entries via their product relation.
The connected cover, coefficient local system, primitive
rational group-algebra projector, actual six-dimensional
Jacobian corner and restricted theta polarization are
unchanged. Their cup product, Poincare biextension,
variational de Rham and Adams realizations carry the
identity maps, with polarization multiplier one, as in
F01--F04. No separate multiplicity representation or
centralizer-quotient cover is used as a detector source.
The proof consists of the actual finite-group averaging,
evaluation, conjugation, rank and quotient maps above;
no classification or unrefereed result supplies a premise.
